A diferença observada entre os valores de \gls{rmse} simulado e experimental
pode ser atribuída a fatores não considerados pelo modelo \gls{fopdt} utilizado
no projeto. Entre as principais causas estão o ruído de medição dos sensores, os
atrasos de transporte e a dinâmica térmica não modelada da planta física.

Outra fonte de discrepância decorre da própria estrutura do sistema, onde o
controlador utiliza a medição do sensor 1 para a realimentação, enquanto o erro
é calculado com base na temperatura medida pelo sensor 2. A diferença de
localização entre os sensores introduz um atraso entre a variável controlada e a
variável de interesse, o que não foi totalmente representado no modelo utilizado no projeto.

A perturbação introduzida aos 10 minutos, decorrente do acionamento do
ventilador, também contribuiu para o aumento do erro experimental, uma vez que
seu efeito real sobre a planta difere da representação simplificada adotada na
simulação.

Apesar dessas divergências, a proximidade entre os valores de
\gls{rmse} confirma que o modelo identificado representa adequadamente o
comportamento da planta e que a estratégia de controle proposta mantém bom
desempenho no seguimento da referência, também em condições experimentais.
